\chapter{Methods and Implemented Components}

\section{Dataset and fitness evaluation}

The common empirical object used in the F1 part of the thesis is a dataset of Framsticks F1 genotype strings paired with scalar fitness values. Each record consists of a genotype written as text and a corresponding \texttt{vertpos} value separated from it by a tab character. The dataset contains 10,893 non-empty records, and all records provide a finite numeric fitness value. This format preserves a direct link between the symbolic description of an organism and the value obtained for that organism in a previous Framsticks evaluation.

In the implemented data-loading procedure, the genotype string and the fitness label are parsed as two separate fields. The genotype remains the primary structured object: it is later encoded either as a character sequence or through grammar-derived production rules, depending on the model family. The fitness value is stored as a floating-point label associated with the same genotype. For character-level processing, internal whitespace is preserved as part of the input string. For grammar-based processing and benchmark seed preparation, whitespace inside the genotype is normalized away so that the genotype can be parsed and compared consistently. This distinction reflects the fact that the dataset is used both for learning representations of genotype strings and for selecting or analysing examples according to their evaluated fitness.

The stored \texttt{vertpos} values should not be interpreted as properties directly computed from raw genotype text. A genotype first has to be accepted by the Framsticks representation, developed into a phenotype, and evaluated in simulation. The \texttt{vertpos} criterion is therefore a phenotype-level measurement returned by Framsticks for the organism described by the genotype. This distinction is important for both parts of the thesis. CoSO and Latent Space Optimization operate on genotype representations or propose genotype-level modifications, but their optimization relevance is determined only after the resulting candidates are interpreted and evaluated as Framsticks organisms.

For autoencoder training, the dataset provides examples from which the syntax and structural regularities of F1 genotypes can be learned. The implementation supports an optional train--validation split rather than relying on a separate validation file. When this split is enabled, a fixed fraction of the parsed examples is assigned to validation using a pseudorandom split controlled by the training seed, while the remaining examples form the training subset. The validation-enabled training scripts use a conventional 0.2 validation fraction. Fitness-aware variants additionally require finite stored fitness labels, because the scalar value can be used as an auxiliary target or as a conditioning signal, but the label remains tied to the original evaluated genotype.

During optimization experiments, dataset fitness values and fresh fitness evaluations play different roles. Stored values are suitable for selecting starting genotypes, grouping examples by initial fitness, and providing labels during representation learning or analysis. However, newly produced candidates are not assigned fitness by looking up or estimating their dataset labels. Instead, candidate genotypes generated by direct genotype-level search, CoSO-style modification, or decoding from latent vectors must be evaluated through Framsticks with the same \texttt{vertpos} criterion. If a candidate cannot be evaluated successfully, it is treated as invalid for optimization purposes. This procedure keeps the final objective common across the compared methods: regardless of how a candidate is produced, its useful fitness is the true Framsticks evaluation of the phenotype generated from its F1 genotype.

\section{Grammar-based representation of F1 genotypes}

The grammar-based part of the implementation represents a Framsticks F1 genotype as a structured expression rather than only as an unconstrained sequence of characters. The raw genotype is still a variable-length string, but for grammar-aware processing it is first normalized by removing internal whitespace and then tokenized into individual F1 symbols. The resulting token sequence is accepted only if it can be parsed by the F1 context-free grammar used in the implementation. This parsing step separates syntactically admissible training examples from strings that cannot be represented as grammar derivations within the configured maximum production-sequence length.

The grammar starts from the nonterminal \texttt{S} and expands it into a sequence of elements. Each element consists of zero or more modifiers followed by a core expression. The core may be the terminal \texttt{X}, which denotes the basic F1 structural symbol used in this simplified grammar, or a parenthesized branch expression. Branches are represented with parentheses and comma-separated branch tails, and several grammar rules allow empty expansions. Empty expansions are important because they make optional components expressible: a modifier sequence may end, a branch list may be absent, and a branch tail may terminate. The modifier nonterminal expands to the set of supported F1 modifier symbols, including both uppercase and lowercase variants such as \texttt{R}/\texttt{r}, \texttt{Q}/\texttt{q}, \texttt{C}/\texttt{c}, \texttt{L}/\texttt{l}, \texttt{W}/\texttt{w}, \texttt{M}/\texttt{m}, \texttt{I}/\texttt{i}, \texttt{F}/\texttt{f}, \texttt{A}/\texttt{a}, \texttt{S}/\texttt{s}, and \texttt{E}/\texttt{e}. At this methodological level, the grammar therefore captures the subset of F1 syntax needed by the grammar-aware models: linear sequences, optional modifier prefixes, nested parenthesized branches, branch separators, and terminal construction symbols.

After parsing, a genotype is encoded as the sequence of productions used in one parse tree. Each production is assigned an integer index, and the production sequence is padded to a fixed maximum length for batch processing. In the integer production-index encoding, positions beyond the end of the actual derivation use a separate padding index; this index is not one of the grammar productions. The one-hot grammar encoding uses a different convention: its columns correspond only to grammar productions, and padded positions activate the final production column, which corresponds to the grammar's \texttt{Nothing -> None} production. These two related dataset encodings support different model interfaces. The index form is convenient for embedding-based sequence models, whereas the one-hot form is convenient for models that predict a distribution over grammar productions. In both cases, the object being learned is not a character string directly, but a derivation trace that can be expanded back into a genotype string.

Decoding follows the inverse idea. A stack of nonterminals is initialized with the start symbol. At each step, the next predicted production is applied to the current nonterminal, terminal symbols are emitted to the reconstructed genotype, and newly introduced nonterminals are pushed onto the stack in the order required for subsequent expansion. This stack-based procedure gives the production sequence an explicit structural interpretation. It also explains why the order of production choices matters: the same production index can be meaningful only when the nonterminal at the top of the derivation stack has a compatible left-hand side.

The implementation therefore precomputes grammar masks from the set of productions. For every left-hand-side nonterminal, a binary mask identifies exactly the productions whose left-hand side matches that nonterminal. During masked decoding, the current derivation stack determines which nonterminal is being expanded, and the corresponding mask suppresses productions that would be syntactically incompatible at that position. When the stack is empty, the remaining sequence positions are treated as a finished derivation through the designated empty or padding state used by the corresponding encoding. This mechanism does not evaluate the biological or fitness quality of a genotype, but it prevents a decoder from selecting productions that violate the local grammar state.

Several later model variants use the same grammar representation with different amounts of structural context. Grammar-rule models operate on production sequences. Masked variants additionally use the current nonterminal to restrict admissible production choices during decoding. Tree-oriented variants reconstruct the left-hand-side context implied by the derivation stack and can provide this context to the encoder or decoder. The LHS-conditioned variants make this context explicit by embedding the currently expanded nonterminal, while LHS-depth variants also represent the current stack depth. These details belong to the autoencoder section at the architectural level, but at the representation level they all rely on the same principle: an F1 genotype is treated as a grammar derivation with a current expansion state, rather than as an unstructured string.

This representation is useful for Latent Space Optimization because latent search ultimately proposes candidates through a decoder. If the decoder is allowed to generate arbitrary symbols, many latent points can map to strings that do not respect the F1 syntax and therefore cannot be meaningfully evaluated. A production-based representation narrows the decoding problem to choosing grammar rules, and grammar masks further constrain these choices to the rules allowed by the current derivation context. Syntactic validity is not sufficient for high \texttt{vertpos} fitness, and it does not guarantee successful optimization, but it is a prerequisite for using decoded candidates as Framsticks genotypes in later evaluation and search procedures.

\section{Tree-based fragment handling}

The CoSO component operates on textual Framsticks F1 genotypes, but substitutions are not selected by cutting the string at arbitrary character positions. Before fragments are considered for replacement, the genotype is interpreted as a tree-like structure induced by the F1 notation. This representation reflects the fact that an F1 genotype contains nested branches and sequences of construction elements, so many syntactically meaningful fragments correspond to complete substructures rather than to unconstrained substrings.

The purpose of this tree-based handling is to define candidate cut points that are likely to preserve the syntactic organization of the genotype. A substitution performed at such boundaries replaces a coherent fragment, for example a branch or a sequence of elements, instead of splitting a modifier sequence, parenthesized branch, or separator structure in the middle. This does not make every substitution valid, because the surrounding context may still require additional syntactic adjustments, but it removes most of the errors that would result from purely character-level cutting. In practice, the resulting candidate genotype is still passed through the Framsticks repair mechanism before evaluation, so minor syntactic inconsistencies such as missing separators or unmatched structural markers may be corrected when possible.

The same tree-based view also provides the pool of fragments that may later be reused as substitutions. Each parsed genotype contributes a set of coherent substructures that can be compared with fragments from other genotypes. Therefore, the fragment extraction stage plays two roles: it identifies the locations in a recipient genotype where a substitution may be attempted, and it supplies donor fragments that can be stored and reused when a compatible context is encountered later in the run.

\section{Compatible Substitutions Optimization module}

The implemented CoSO module performs optimization directly in the F1 genotype space. It does not introduce an intermediate continuous representation and does not train a model before the search. Instead, it maintains a collection of previously observed genotype fragments and uses them to propose context-aware substitutions in currently processed genotypes. The final objective remains the true Framsticks evaluation of the decoded organism with the \texttt{vertpos} criterion.

At the beginning of a run, the module initializes the Framsticks evaluation environment, random seeds, population parameters, archive parameters, compatibility thresholds, and the logging procedure. Starting genotypes can be obtained either from a dataset of F1 genotypes with stored fitness values or generated through the Framsticks mutation mechanism. Dataset fitness values are used when the exact genotype is already known, whereas newly produced candidates are evaluated through Framsticks. Empty genotypes and genotypes exceeding the configured length constraint are treated as uncompetitive, and all accepted evaluations update the best solution observed so far.

The search is organized around an active genotype that is repeatedly improved by compatible substitutions. For this genotype, the tree-based fragment handling described above produces candidate recipient fragments together with their locations in the original genotype. Each fragment is described by a small set of representation-level features. In the current F1 implementation, these features characterize the spatial effect of the fragment and the direction associated with its terminal part. The intent is to compare fragments by the structural role they may play in the constructed phenotype, rather than by their raw textual form or absolute position in the genotype.

The archive stores donor fragments indexed by these feature values. When a genotype can no longer be improved, its fragments are inserted into the archive together with the fitness of the solution from which they were obtained. For each feature bucket, only a limited number of entries is retained, with preference given to fragments originating from better-performing genotypes. The archive therefore acts as a memory of building blocks discovered during the search, grouped by approximate compatibility rather than by genotype position.

During improvement, the module queries the archive for fragments whose feature values are close enough to the feature values of a recipient fragment. The accepted distance is controlled by the spatial and directional compatibility thresholds. Retrieved donor fragments are then tested as possible replacements. Before a candidate is evaluated, overly large substitutions and substitutions that would violate the maximum genotype length are rejected. This prevents a compatible feature match from replacing too much of the current genotype or producing candidates outside the intended representation constraints.

For each remaining candidate, the selected recipient fragment is replaced with the donor fragment in the full genotype. The candidate is then passed to the Framsticks repair mechanism. This repair step is important because feature compatibility and tree-based cut points reduce syntactic disruption but do not guarantee that every resulting string is directly accepted by Framsticks. Candidates that cannot be repaired or evaluated are discarded. Candidates that can be evaluated are accepted only if their \texttt{vertpos} value is strictly higher than the fitness of the genotype before the substitution. Thus, as in optimal mixing, a fragment is considered useful only after it has been tested in the complete recipient context.

If an accepted substitution is found, the improved genotype becomes the new active genotype and the process continues from this modified solution. If no compatible substitution improves the active genotype, the solution is treated as locally converged. Its fragments are added to the archive, the solution may be inserted into the elite set, and a new active genotype is selected or generated. The elite set is kept limited to the best solutions encountered during such convergence events. This creates a restart-like process in which later genotypes are improved using an archive that accumulates fragments from earlier search attempts.

The module records the main events of the run in a JSONL log. The stored events include metadata, new starting genotypes, accepted improvements, archive retrieval statistics, convergence events, progress snapshots, final elite solutions, and the best genotype found. These logs are used later to reconstruct optimization curves and analyse how the archive-based substitution process behaves over time. The run terminates when the configured number of full Framsticks evaluations is reached.

\section{Autoencoder-based genotype representation}

The autoencoder-based representation used in the Latent Space Optimization part of the thesis maps a variable-length F1 genotype to a fixed-dimensional numerical vector and then reconstructs a genotype from that vector. In all variants, the encoder is responsible for transforming an input sequence into a latent representation, while the decoder is responsible for producing a sequence that can be interpreted as a genotype. This representation is not a replacement for the Framsticks genotype itself. It is an intermediate search representation: optimization may operate on the latent vector, but every candidate must eventually be decoded back to F1 text before it can be evaluated with the true \texttt{vertpos} criterion.

Two input families are used. The first family treats the genotype as a character sequence. A character-level dataset builds a vocabulary containing padding, start, and end symbols, followed by the characters observed in the training strings. Each genotype is encoded as a padded sequence of token indices. The corresponding character-level variational autoencoder uses recurrent sequence encoders and decoders: the encoder embeds the tokens and summarizes the sequence into the parameters of a latent distribution, and the decoder generates output tokens autoregressively from a sampled latent vector. This variant leaves syntactic regularities to be learned from examples, because it does not explicitly know which character choices are legal in a given grammatical context.

The second family uses the grammar representation introduced in Section~3.2. A parsed genotype is encoded as a sequence of production-rule indices, optionally represented in one-hot form. The reconstruction target is therefore not directly a sequence of raw characters, but a derivation trace that can be expanded into an F1 string. This distinction is important for validity. If the decoder predicts a rule sequence that follows the derivation stack consistently, the resulting output is constrained by the same grammar that was used to encode the training examples. The decoder output is then converted to genotype text by applying predicted productions to the current nonterminal and emitting terminal symbols.

The implemented grammar-rule variational autoencoder follows the standard encoder--decoder pattern at the level of production indices. The encoder embeds the production sequence and maps it to mean and log-variance vectors. During training, a latent vector is sampled by the reparameterization mechanism, and the decoder is trained to reproduce the input production sequence. The objective combines a reconstruction term over sequence positions with a Kullback--Leibler regularization term, with a scheduled weight on the regularization term during training. Teacher forcing is used for recurrent decoders during training and is reduced over the training schedule. At the methods level, the relevant point is that the model learns a continuous latent code from which complete grammar-rule sequences can be generated.

Masked grammar variants add explicit syntactic constraints to decoding. The F1 grammar provides, for each left-hand-side nonterminal, the set of productions that may expand it. During masked decoding, the current derivation stack determines the active nonterminal, and logits for incompatible productions are suppressed. This mechanism does not guarantee that a decoded genotype has high fitness, and it does not evaluate semantic properties of the organism. It only enforces local grammatical compatibility during rule selection, which reduces the number of syntactically impossible production choices available to the decoder.

Tree-oriented variants use the fact that a production sequence implicitly defines a derivation tree. In these models, the current left-hand-side context is reconstructed from the stack of nonterminals generated by previous rules and is provided as an additional input to the encoder. Masked tree variants combine this tree-aware context with rule masks in the decoder. LHS-conditioned masked tree variants go further by giving the decoder an embedding of the currently expanded nonterminal, and LHS-depth variants additionally represent the current stack depth. These mechanisms all use the same grammar as Section~3.2, but they expose more of the derivation state to the neural model instead of requiring it to infer the state only from preceding rule indices.

Conditional extensions are also implemented for selected tree-oriented models. Fitness-conditioned variants append a scalar conditioning value to the latent vector before decoding. Structural conditional variants use a small vector of auxiliary structural descriptors rather than only a scalar fitness condition. These variants are methodological extensions of the same autoencoder design: the latent vector remains the main representation, while the condition supplies additional information to the decoder. The presence of a conditioning signal should not be interpreted as replacing true fitness evaluation, because decoded candidates are still evaluated through Framsticks during optimization.

A Transformer-based grammar variational autoencoder is included as an alternative to recurrent sequence modelling. It uses embeddings of production-rule indices together with positional encodings and Transformer encoder--decoder blocks. The encoder pools non-padding positions to obtain latent distribution parameters, and the decoder generates a production sequence from the latent vector using a causal decoding structure. This variant keeps the same grammar-rule input and reconstruction target as other grammar-based models, but changes the sequence model used to capture dependencies between production choices.

Finally, a vector-quantized grammar autoencoder is implemented as a related representation-learning variant. Instead of sampling from a continuous variational distribution, it maps an encoded production sequence to a latent vector that is quantized through a learned codebook and then decoded back into production rules. This variant is relevant as an alternative autoencoding mechanism for grammar-rule sequences, although the main Latent Space Optimization procedures considered in this thesis operate on continuous latent vectors generated by variational models.

Across these model families, the training target remains reconstruction of the input representation: characters for the character-level model and grammar production sequences for grammar-aware models. Reconstruction, however, is only a prerequisite for later optimization. A model that reconstructs valid examples well may still define a latent space that is difficult to optimize, and a syntactically constrained decoder may still produce organisms with low \texttt{vertpos}. For this reason, autoencoder training and latent-space search are treated as separate methodological components: the autoencoder defines the representation, while the optimization methods described next decide how that representation is explored.

\section{Latent-space optimization methods}

Latent-space optimization uses a trained genotype autoencoder as a generator of candidate F1 genotypes. The optimizer itself manipulates fixed-dimensional latent vectors, but the objective value assigned to a vector is obtained only after decoding. In the implemented latent EA, CEM, and CMA-ES benchmark procedures, latent candidates are decoded as grammar-rule sequences, converted into genotype text, checked for syntactic admissibility, and evaluated in Framsticks. Character-level autoencoders remain part of the broader representation set described in Section~3.5, but they are not the decoding context used by these grammar-rule optimization loops. The scalar feedback used by the optimizer is the resulting true \texttt{vertpos} fitness, with invalid or unevaluable candidates treated as uncompetitive. Thus, the learned model proposes candidates, whereas Framsticks remains the source of fitness values.

In the seeded setting used by the implemented procedures, optimization starts from one or more existing genotypes. A starting genotype is encoded by the autoencoder, and the mean vector produced by the encoder is used as the initial latent point. This choice gives the search a concrete reference candidate rather than beginning only from an unconditional prior sample. For grammar-rule models, starting genotypes must be encodable by the grammar representation and must fit within the configured maximum production-sequence length. The decoded form of the starting latent vector can also be evaluated, because the autoencoder reconstruction of the seed may differ from the original genotype.

All latent-space methods share the same evaluation loop. A population or distribution produces a batch of latent vectors. Each vector is decoded to an F1 genotype, normalized for comparison, and rejected if it is empty, syntactically invalid, duplicated when uniqueness is required, or impossible to evaluate. Remaining candidates are evaluated by Framsticks with the \texttt{vertpos} criterion. Fitness values are cached for already evaluated genotypes so that the optimizer is not rewarded for repeatedly generating the same string. The best decoded genotype and its true fitness are tracked separately from the latent vector that produced it, because the final output of the method is an F1 genotype rather than the vector itself.

The latent-space evolutionary search maintains a population of individuals represented by both decoded genotype and latent vector. Initial population members are produced by perturbing the encoded seed vector with Gaussian noise of selected standard deviations and evaluating the decoded genotypes. In each generation, parents are selected according to true evaluated fitness. Offspring latent vectors are generated through several continuous-space operators: mutation by Gaussian perturbation, crossover-like interpolation between parent vectors, optional extrapolation along lines between parents, and directional moves biased by the difference between a high-fitness vector and another selected vector. After decoding and true evaluation, valid offspring are merged with elite survivors, and the population is truncated to the configured size.

Random immigrants are supported in the latent evolutionary search as an exploration mechanism. With a configurable probability, a child vector is sampled from a standard normal distribution rather than produced from selected parents. This operation introduces latent candidates that are not constrained to remain near the current population. Methodologically, it reduces the risk that the population collapses prematurely around a narrow decoded region, although the sampled vectors must still decode to valid and useful genotypes to affect the search.

The same evolutionary implementation can optionally re-encode accepted offspring. In this variant, a decoded genotype that passes validity and evaluation is encoded again, and the resulting latent mean replaces the raw child vector. Re-encoding projects the individual back through the autoencoder's representation of the decoded genotype. This can make the stored latent population more consistent with the model's encoder--decoder manifold, but it is an operator choice rather than a separate objective function.

The Cross-Entropy Method (CEM) is implemented as a distribution-based optimizer over latent vectors. Starting from the encoded seed, the method maintains a diagonal normal sampling distribution. At each iteration, it samples a population of latent candidates, decodes and evaluates them, selects an elite fraction according to \texttt{vertpos}, and updates the distribution mean and standard deviation toward the elite set. Smoothing and a minimum standard deviation are used so that the distribution does not change too abruptly or collapse completely. CEM therefore differs from the latent evolutionary search in that it does not maintain explicit genotype individuals across generations; it adapts a sampling distribution from elite latent samples.

CMA-ES is implemented as another continuous distribution-search method. It samples latent candidates around a mean vector, evaluates the decoded genotypes, and updates the search mean and covariance structure according to the best-performing samples. In this setting, the optimizer minimizes the negative of the true fitness only implicitly through the scores returned after decoding and Framsticks evaluation; it does not use gradients through the decoder or through the simulator.

The latent-space methods are therefore derivative-free black-box optimizers with respect to the final objective. The autoencoder is trained before optimization and is used only to map between F1 genotypes and latent vectors. During search, no learned surrogate fitness is used to replace Framsticks evaluation, and no gradient of \texttt{vertpos} is available. This separation is essential for the comparison performed in later chapters: differences between methods can be attributed to how they explore the learned representation and how often they produce valid, diverse, high-fitness decoded genotypes, not to a change in the underlying fitness criterion.

\section{Baseline evolutionary methods}

% TODO: write this section.

\section{Implementation architecture}

% TODO: write this section.
